%%%%%%%%%%%%%%%%%%%%%%%%%%%%%%%%%%%%%%%%%%%%%%%%%%%%%%%%%%%%%%%%%%%%%%%%%%%%%%%
%
%   SZABLON SPRAWOZDANIA - AOD (JĘZYK JULIA / JUMP)
%   Autor: Marcel Musiałek, 279704
%
%   Na podstawie szablonu z Obliczeń Naukowych
%
%%%%%%%%%%%%%%%%%%%%%%%%%%%%%%%%%%%%%%%%%%%%%%%%%%%%%%%%%%%%%%%%%%%%%%%%%%%%%%%

% --- Klasa dokumentu i podstawowe ustawienia ---
\documentclass[11pt, a4paper]{article}

% --- Podstawowe pakiety kodowania i języka ---
\usepackage[utf8]{inputenc} % Kodowanie wejściowe
\usepackage[T1]{fontenc}    % Kodowanie fontów (obsługa polskich znaków)
\usepackage{lmodern}      % Nowoczesna czcionka Latin Modern
\usepackage[polish]{babel}  % Ustawia polskie nazwy (np. "Spis treści", "Rysunek")

% --- Pakiety matematyczne ---
\usepackage{amsmath} % Zaawansowane formuły matematyczne
\usepackage{amssymb} % Dodatkowe symbole, np. \forall

% --- Ustawienia strony i układu ---
\usepackage{geometry}
\geometry{a4paper, margin=2.5cm} % Marginesy
\usepackage{verbatim}  % Dla środowiska \begin{verbatim} (surowy tekst)
\usepackage{enumitem}  % Lepsza kontrola nad listami

% --- Pakiet Hyperref (dla linków i zakładek PDF) ---
\usepackage[unicode=true, pdfusetitle]{hyperref}
\hypersetup{
    colorlinks=true,
    linkcolor=blue,
    filecolor=magenta,      
    urlcolor=cyan,
    pdftitle={Sprawozdanie AOD},
}

%%%%%%%%%%%%%%%%%%%%%%%%%%%%%%%%%%%%%%%%%%%%%%%%%%%%%%%%%%%%%%%%%%%%%%%%%%%%%%%
%
%   STRONA TYTUŁOWA DOKUMENTU
%
%%%%%%%%%%%%%%%%%%%%%%%%%%%%%%%%%%%%%%%%%%%%%%%%%%%%%%%%%%%%%%%%%%%%%%%%%%%%%%%

\title{Algorytmy Optymalizacji Dyskretnej \\ Laboratorium 2: Programowanie Liniowe}
\author{Marcel Musiałek \\ 279704}
\date{\today}


%%%%%%%%%%%%%%%%%%%%%%%%%%%%%%%%%%%%%%%%%%%%%%%%%%%%%%%%%%%%%%%%%%%%%%%%%%%%%%%
%
%   GŁÓWNY DOKUMENT
%
%%%%%%%%%%%%%%%%%%%%%%%%%%%%%%%%%%%%%%%%%%%%%%%%%%%%%%%%%%%%%%%%%%%%%%%%%%%%%%%

\begin{document}

% Wygeneruj stronę tytułową
\maketitle

% Spis treści
\tableofcontents
\newpage

%%%%%%%%%%%%%%%%%%%%%%%%%%%%%%%%%%%%%%%%%%%%%%%%%%%%%%%%%%%%%%%%%%%%%%%%%%%%%%%
%
%   TREŚĆ SPRAWOZDANIA
%
%%%%%%%%%%%%%%%%%%%%%%%%%%%%%%%%%%%%%%%%%%%%%%%%%%%%%%%%%%%%%%%%%%%%%%%%%%%%%%%

% --- POCZĄTEK ZADANIA 1 ---
\section{Zadanie 1: Problem dostaw paliwa}

\subsection{Opis Modelu}
Zmienne i parametry modelu:

\begin{itemize}
    \item $n$ - ilość lotnisk ($n=4$)
    \item $m$ - ilość dostawców paliwa ($m=3$)
    \item $I = \{1, \dots, n\}$ - zbiór lotnisk
    \item $J = \{1, \dots, m\}$ - zbiór dostawców
    \item $d \in \mathbb{R}^n$ - wektor popytu (demand), zapotrzebowanie na paliwo na $i$-tym lotnisku.
    \item $s \in \mathbb{R}^m$ - wektor podaży (supply), maksymalna ilość paliwa od $j$-tego dostawcy.
    \item $C = (c_{ij}) \in \mathbb{R}^{n \times m}$ - macierz kosztów, cena [\$ / galon] dostarczenia paliwa na $i$-te lotnisko przez $j$-tego dostawcę.
\end{itemize}

\noindent \textbf{Zmienne decyzyjne:}
$$ X = (x_{ij}) \in \mathbb{R}^{n \times m} $$
gdzie $x_{ij}$ to ilość paliwa [w galonach] dostarczona przez $j$-tego dostawcę na $i$-te lotnisko.

\noindent \textbf{Cel - minimalizacja kosztów:}
$$ \min \sum_{i=1}^{n} \sum_{j=1}^{m} c_{ij} \cdot x_{ij} $$

\noindent \textbf{Ograniczenia:}
\begin{align*}
    x_{ij} &\ge 0 & \forall i \in I, \forall j \in J \quad &\text{(Nieujemność dostaw)} \\
    \sum_{j=1}^{m} x_{ij} &= d_i & \forall i \in I \quad &\text{(Zaspokojenie popytu)} \\
    \sum_{i=1}^{n} x_{ij} &\le s_j & \forall j \in J \quad &\text{(Limit podaży)}
\end{align*}

\subsection{Egzemplarz z listy}
$n = 4$, $m = 3$
$$ d =
\begin{pmatrix}
110000 \\
220000 \\
330000 \\
440000
\end{pmatrix}, \quad s =
\begin{pmatrix}
275000 \\
550000 \\
660000
\end{pmatrix}, \quad C =
\begin{pmatrix}
10 & 7 & 8 \\
10 & 11 & 14 \\
9 & 12 & 4 \\
11 & 13 & 9
\end{pmatrix}
$$
\bigskip
\noindent \textbf{Rozwiązanie (za pomocą GLPK):}
Optymalna macierz dostaw $X$ (wiersze = lotniska, kolumny = firmy):
$$ X =
\begin{pmatrix}
0 & 110000 & 0 \\
165000 & 55000 & 0 \\
0 & 0 & 330000 \\
110000 & 0 & 330000
\end{pmatrix}
$$

\noindent \textbf{Interpretacja wyników:}
\begin{itemize}
    \item[ (a)] \textbf{Minimalny łączny koszt} dostaw wynosi \textbf{8 525 000 \$}.
    \item[ (b)] \textbf{Wszystkie firmy dostarczają paliwo} (każda kolumna macierzy $X$ zawiera co najmniej jedną niezerową wartość).
    \item[ (c)] \textbf{Możliwości Firmy 1 i Firmy 3 są wyczerpane} (sumy kolumn 1 i 3 są równe $s_1$ i $s_3$). Firma 2 dostarcza łącznie 165 000 galonów, wykorzystując 30\% swoich możliwości (165 000 / 550 000).
\end{itemize}

% --- KONIEC ZADANIA 1 ---

% --- POCZĄTEK ZADANIA 2 ---
\section{Zadanie 2: Optymalny plan produkcji}

\subsection{Opis Modelu}
Problem polega na maksymalizacji tygodniowego zysku z produkcji czterech wyrobów, przy ograniczeniach czasowych trzech maszyn oraz ograniczeniach popytu.

\begin{itemize}
    \item $W = \{P1, P2, P3, P4\}$ - zbiór wyrobów ($i \in W$)
    \item $M = \{M1, M2, M3\}$ - zbiór maszyn ($j \in M$)
    \item $t_{ij}$ - czas obróbki [min/kg] wyrobu $i$ na maszynie $j$
    \item $d_i$ - maksymalny tygodniowy popyt [kg] na wyrób $i$
    \item $a_j$ - dostępny czas pracy [min] maszyny $j$ (3600 min dla każdej)
    \item $z_i$ - zysk jednostkowy [\$/kg] dla wyrobu $i$. Jest to wartość obliczona jako:
    \[ z_i = (\text{Cena}_i) - (\text{KosztMat}_i) - \sum_{j \in M} (t_{ij} \cdot \text{KosztPracyMin}_j) \]
\end{itemize}

\noindent \textbf{Zmienne decyzyjne:}
$$ x_i \quad \text{dla } i \in W $$
gdzie $x_i$ to tygodniowa produkcja [w kg] wyrobu $i$.

\noindent \textbf{Cel - maksymalizacja zysku:}
$$ \max \sum_{i \in W} z_i \cdot x_i $$

\noindent \textbf{Ograniczenia:}
\begin{align*}
    \sum_{i \in W} t_{ij} \cdot x_i &\le a_j & \forall j \in M \quad &\text{(Limit czasu maszyn)} \\
    0 \le x_i &\le d_i & \forall i \in W \quad &\text{(Limit popytu i nieujemność)}
\end{align*}

\subsection{Egzemplarz z listy}
Dane wejściowe (z tabel i treści zadania):
$$ t =
\begin{pmatrix}
5 & 10 & 6 \\
3 & 6 & 4 \\
4 & 5 & 3 \\
4 & 2 & 1
\end{pmatrix} 
\quad d =
\begin{pmatrix}
400 \\ 100 \\ 150 \\ 500
\end{pmatrix} 
\quad a =
\begin{pmatrix}
3600 \\ 3600 \\ 3600
\end{pmatrix} 
$$
Obliczony wektor zysku jednostkowego (po odjęciu kosztów materiałowych i kosztów pracy maszyn):
$$ z =
\begin{pmatrix}
4.200 \\ 5.500 \\ 4.550 \\ 3.750
\end{pmatrix}
$$

\bigskip
\noindent \textbf{Rozwiązanie (za pomocą GLPK):}
Optymalny wektor produkcji $x$:
$$ x =
\begin{pmatrix}
125.00 \\
100.00 \\
150.00 \\
500.00
\end{pmatrix}
$$

\noindent \textbf{Interpretacja wyników:}
\begin{itemize}
    \item Optymalny tygodniowy plan produkcji zakłada produkcję 125 kg P1, 100 kg P2, 150 kg P3 oraz 500 kg P4.
    \item Produkcja wyrobów P2, P3 i P4 jest \textbf{ograniczona popytem} (osiągnęła 100\% maksymalnej sprzedaży). Produkcja P1 jest ograniczona przez zasoby.
    \item Maszyna M2 pracuje z pełną wydajnością (100\%), M1 jest bliska limitu (97.92\%), a M3 ma duży zapas mocy (58.33\%). Oznacza to, że \textbf{maszyna M2 jest "wąskim gardłem"} procesu.
    \item Osiągnięty \textbf{maksymalny zysk tygodniowy} wynosi \textbf{3632.50 \$}.
\end{itemize}

% --- KONIEC ZADANIA 2 ---

% --- POCZĄTEK ZADANIA 3 ---
\section{Zadanie 3: Planowanie produkcji i magazynowania}

\subsection{Opis Modelu}
Problem wielookresowego planowania produkcji z możliwością magazynowania i produkcją ponadwymiarową.

\begin{itemize}
    \item $K$ - liczba okresów ($K=4$)
    \item $J = \{1, \dots, K\}$ - zbiór okresów
    \item $d_j$ - popyt na towar w okresie $j \in J$
    \item $c_j$ - koszt jednostkowy produkcji normalnej w okresie $j \in J$
    \item $o_j$ - koszt jednostkowy produkcji ponadwymiarowej w okresie $j \in J$
    \item $a_j$ - limit produkcji ponadwymiarowej w okresie $j \in J$
    \item $h$ - koszt jednostkowy magazynowania (stały = 1500)
    \item $P_{\max}$ - limit produkcji normalnej (stały = 100)
    \item $M_{\max}$ - limit pojemności magazynu (stały = 70)
    \item $m_0$ - stan magazynu na początku (stały = 15)
\end{itemize}

\noindent \textbf{Zmienne decyzyjne:}
\begin{itemize}
    \item $p_j$ - ilość produkcji normalnej w okresie $j \in J$
    \item $e_j$ - ilość produkcji ponadwymiarowej w okresie $j \in J$
    \item $m_j$ - stan magazynu na koniec okresu $j \in J \cup \{0\}$
\end{itemize}

\noindent \textbf{Cel - minimalizacja kosztów:}
$$ \min \left( \sum_{j=1}^{K} (c_j \cdot p_j + o_j \cdot e_j) + \sum_{j=1}^{K} (h \cdot m_j) \right) $$

\noindent \textbf{Ograniczenia:}
\begin{align*}
    m_j &= m_{j-1} + p_j + e_j - d_j & \forall j \in J \quad &\text{(Bilans magazynu)} \\
    p_j &\le P_{\max} & \forall j \in J \quad &\text{(Limit prod. normalnej)} \\
    e_j &\le a_j & \forall j \in J \quad &\text{(Limit prod. ponadwymiarowej)} \\
    m_j &\le M_{\max} & \forall j \in J \quad &\text{(Limit magazynu)} \\
    m_0 &= 15 & \quad &\text{(Warunek początkowy)} \\
    p_j, e_j, m_j &\ge 0 & \forall j \in J &
\end{align*}

\subsection{Egzemplarz z listy}
Dane wejściowe (z tabeli):
$$ d = \begin{pmatrix} 130 \\ 80 \\ 125 \\ 195 \end{pmatrix} \quad
   c = \begin{pmatrix} 6000 \\ 4000 \\ 8000 \\ 9000 \end{pmatrix} \quad
   o = \begin{pmatrix} 8000 \\ 6000 \\ 10000 \\ 11000 \end{pmatrix} \quad
   a = \begin{pmatrix} 60 \\ 65 \\ 70 \\ 60 \end{pmatrix}
$$
Oraz $K=4, P_{\max}=100, M_{\max}=70, h=1500, m_0=15$.

\bigskip
\noindent \textbf{Rozwiązanie (za pomocą GLPK):}

Wektory zmiennych decyzyjnych w rozwiązaniu optymalnym:
$$ p = \begin{pmatrix} 100 \\ 100 \\ 100 \\ 100 \end{pmatrix} \quad
   e = \begin{pmatrix} 15 \\ 50 \\ 0 \\ 50 \end{pmatrix} \quad
   m = \begin{pmatrix} 0 \\ 70 \\ 45 \\ 0 \end{pmatrix}
$$
(Wektor $m$ pokazuje stan na koniec okresów 1, 2, 3, 4; stan $m_0=15$).

\noindent \textbf{Interpretacja wyników:}
\begin{itemize}
    \item[ (a)] \textbf{Minimalny łączny koszt} produkcji i magazynowania wynosi \textbf{3 842 500 \$}.
    \item[ (b)] Firma musi zaplanować \textbf{produkcję ponadwymiarową w okresach 1, 2 oraz 4} (wektor $e$ ma tam wartości 15, 50, 50). W okresie 3 produkcja ponadwymiarowa nie jest opłacalna.
    \item[ (c)] Możliwości magazynowania (70 jednostek) są \textbf{wyczerpane na koniec okresu 2} (gdzie $m_2 = 70$). Firma celowo produkuje "na zapas" w tanim okresie 2 (koszt prod. norm 4000, ponadwym. 6000), płaci za magazynowanie (1500), aby uniknąć znacznie droższej produkcji w okresie 3 (koszt norm 8000, ponadwym. 10000).
\end{itemize}

% --- KONIEC ZADANIA 3 ---

% --- POCZĄTEK ZADANIA 4 ---
\section{Zadanie 4: Najkrótsza ścieżka z ograniczeniem}

\subsection{Opis Modelu}
Problem znajdowania najkrótszej ścieżki z ograniczeniem zasobu (RCSP). Jest to problem NP-trudny, ale dla małych instancji można go efektywnie rozwiązać za pomocą programowania całkowitoliczbowego (MIP). Model bazuje na sformułowaniu przepływowym.

\begin{itemize}
    \item $N$ - zbiór wierzchołków (miast)
    \item $A \subseteq N \times N$ - zbiór łuków (połączeń), $(i, j) \in A$
    \item $c_{ij}$ - koszt przejazdu łukiem $(i, j)$
    \item $t_{ij}$ - czas przejazdu łukiem $(i, j)$
    \item $i_0 \in N$ - wierzchołek startowy
    \item $j_0 \in N$ - wierzchołek końcowy
    \item $T$ - maksymalny dopuszczalny czas przejazdu
\end{itemize}

\noindent \textbf{Zmienne decyzyjne:}
$$ x_{ij} \in \{0, 1\} \quad \forall (i, j) \in A $$
gdzie $x_{ij} = 1$, jeśli łuk $(i, j)$ należy do wybranej ścieżki, i $x_{ij} = 0$ w przeciwnym razie.

\noindent \textbf{Cel - minimalizacja kosztu:}
$$ \min \sum_{(i, j) \in A} c_{ij} \cdot x_{ij} $$

\noindent \textbf{Ograniczenia:}
\begin{align*}
    \sum_{(i, j) \in A} t_{ij} \cdot x_{ij} &\le T & &\text{(Limit czasu)} \\
    \sum_{(i, k) \in A} x_{ik} - \sum_{(k, i) \in A} x_{ki} &= b_k & \forall k \in N \quad &\text{(Bilans przepływu)}
\end{align*}
gdzie wektor bilansu $b_k$ jest zdefiniowany jako:
$$ b_k = \begin{cases} 
      1 & \text{dla } k = i_0 \text{ (start)} \\
      -1 & \text{dla } k = j_0 \text{ (koniec)} \\
      0 & \text{dla pozostałych } k \in N
   \end{cases}
$$

\subsection{Rozwiązanie (a) - Egzemplarz z listy}
Dane: $N = \{1, \dots, 10\}$, $i_0 = 1$, $j_0 = 10$, $T = 15$. Koszty i czasy zgodnie z listą.

\noindent \textbf{Rozwiązanie (za pomocą GLPK):}
\begin{itemize}
    \item \textbf{Minimalny koszt:} 13 \$
    \item \textbf{Całkowity czas:} 15 (wykorzystano 100\% limitu $T=15$)
    \item \textbf{Optymalna ścieżka:} 1 $\to$ 2 $\to$ 3 $\to$ 5 $\to$ 7 $\to$ 9 $\to$ 10
\end{itemize}

\subsection{Rozwiązanie (b) - Własny egzemplarz}
Stworzono egzemplarz ($n=10, T=15$) z dwiema głównymi ścieżkami:
\begin{itemize}
    \item Ścieżka 1 (2 krawędzie: 1$\to$5$\to$10): Koszt = 2, Czas = 20
    \item Ścieżka 2 (3 krawędzie: 1$\to$2$\to$6$\to$10): Koszt = 6, Czas = 12
\end{itemize}
Najtańsza ścieżka (Koszt 2) jest celowo za długa (20 > 15).

\noindent \textbf{Rozwiązanie (za pomocą GLPK):}
\begin{itemize}
    \item \textbf{Minimalny koszt:} 6 \$
    \item \textbf{Całkowity czas:} 12 (mieści się w limicie $T=15$)
    \item \textbf{Optymalna ścieżka:} 1 $\to$ 2 $\to$ 6 $\to$ 10
\end{itemize}
Model poprawnie odrzucił najtańszą ścieżkę (2-krawędziową) z powodu limitu czasu i znalazł kolejną najlepszą (3-krawędziową), spełniając założenia podpunktu (b).

\subsection{Odpowiedzi na pytania (c) i (d)}

\noindent \textbf{(c) Czy ograniczenie na całkowitoliczbowość jest potrzebne?}

\textbf{Tak, jest absolutnie konieczne.} Model programowania liniowego (LP) bez tego ograniczenia (tj. $x_{ij} \in [0, 1]$) mógłby znaleźć rozwiązanie ułamkowe. Solver LP mógłby "podzielić" przepływ $b_k=1$ na dwie lub więcej ścieżek jednocześnie, aby uśrednić koszty i czasy, co nie reprezentuje fizycznego przejazdu.

\textbf{Kontrprzykład:} Załóżmy, że mamy dwie ścieżki ze startu (1) do celu (2):
\begin{itemize}
    \item Ścieżka A (1$\to$2): Koszt = 10, Czas = 10
    \item Ścieżka B (1$\to$2): Koszt = 0, Czas = 20
    \item Limit czasu $T = 15$.
\end{itemize}
Model MIP poprawnie wybrałby ścieżkę A (Koszt 10). Natomiast model LP (bez całkowitoliczbowości) mógłby ustawić $x_A = 0.5$ i $x_B = 0.5$. Dałoby to "uśredniony" wynik: Całkowity Koszt = $0.5 \cdot 10 + 0.5 \cdot 0 = 5$ oraz Całkowity Czas = $0.5 \cdot 10 + 0.5 \cdot 20 = 15$. Takie rozwiązanie (koszt 5) jest lepsze od rozwiązania całkowitoliczbowego (koszt 10), ale jest bezużyteczne, bo nie reprezentuje żadnej realnej ścieżki.

\noindent \textbf{(d) Czy po usunięciu limitu czasu (w modelu LP) wynik jest akceptowalny?}

\textbf{Tak.} Jeżeli usuniemy ograniczenie na czas (w modelu LP, czyli $x_{ij} \in [0, 1]$) i rozwiążemy problem $\min \sum c_{ij} x_{ij}$ z samymi ograniczeniami bilansu przepływu, otrzymamy klasyczny problem najkrótszej ścieżki (LP). Macierz opisująca ograniczenia bilansu przepływu jest **całkowicie unimodularna**. Oznacza to, że dla dowolnych całkowitoliczbowych danych wejściowych (wektor $b_k$), rozwiązanie optymalne problemu LP (jeśli istnieje) będzie **zawsze całkowitoliczbowe** (tj. $x_{ij}$ będą wynosić 0 lub 1).

W skrócie: dla problemu najkrótszej ścieżki (bez dodatkowych ograniczeń jak czas), nie musimy wymuszać całkowitoliczbowości, ponieważ dostajemy ją "za darmo" z własności matematycznych problemu.

% --- KONIEC ZADANIA 4 ---

% --- POCZĄTEK ZADANIA 5 ---
\section{Zadanie 5: Przydział radiowozów (Problem cyrkulacji)}

\subsection{Opis Modelu}
Problem polega na ustaleniu liczby radiowozów w poszczególnych dzielnicach na poszczególnych zmianach tak, aby spełnić ograniczenia minimalne i maksymalne oraz zminimalizować łączną flotę pojazdów. Zadanie sformułowano jako problem programowania całkowitoliczbowego.

\begin{itemize}
    \item $I = \{p1, p2, p3\}$ - zbiór dzielnic
    \item $J = \{1, 2, 3\}$ - zbiór zmian
    \item $L_{ij}, U_{ij}$ - minimalna i maksymalna liczba radiowozów w dzielnicy $i$ na zmianie $j$ (z tabel 1 i 2)
    \item $S_j$ - wymagana minimalna liczba radiowozów łącznie na zmianie $j$ ($[10, 20, 18]$)
    \item $D_i$ - wymagana minimalna liczba radiowozów łącznie w dzielnicy $i$ ($[10, 14, 13]$)
\end{itemize}

\noindent \textbf{Zmienne decyzyjne:}
$$ x_{ij} \in \mathbb{Z}_{\ge 0} $$
Liczba radiowozów przydzielonych do dzielnicy $i$ na zmianie $j$.

\noindent \textbf{Cel - minimalizacja floty:}
$$ \min \sum_{i \in I} \sum_{j \in J} x_{ij} $$

\noindent \textbf{Ograniczenia:}
\begin{align*}
    L_{ij} \le x_{ij} &\le U_{ij} & \forall i \in I, j \in J \quad &\text{(Limity lokalne)} \\
    \sum_{i \in I} x_{ij} &\ge S_j & \forall j \in J \quad &\text{(Minima dla zmian)} \\
    \sum_{j \in J} x_{ij} &\ge D_i & \forall i \in I \quad &\text{(Minima dla dzielnic)}
\end{align*}

\subsection{Wyniki i interpretacja}
Dla podanych danych solver GLPK znalazł rozwiązanie optymalne.

\noindent \textbf{Optymalny przydział radiowozów (macierz $X$):}
$$
\begin{array}{r|ccc|c}
 & \text{Zmiana 1} & \text{Zmiana 2} & \text{Zmiana 3} & \text{Suma Dzielnicy} \\
\hline
p1 & 2 & 7 & 5 & 14 \\
p2 & 3 & 6 & 7 & 16 \\
p3 & 5 & 7 & 6 & 18 \\
\hline
\text{Suma Zmiany} & 10 & 20 & 18 & \textbf{48 (Całkowita)}
\end{array}
$$

\noindent \textbf{Interpretacja:}
Minimalna łączna liczba radiowozów potrzebna do spełnienia wszystkich wymogów wynosi **48**.
\begin{itemize}
    \item Wymogi dla zmian zostały spełnione dokładnie na poziomie minimów (10, 20, 18).
    \item Wszystkie dzielnice otrzymały liczbę pojazdów przekraczającą wymagane minima ($14 > 10$, $16 > 14$, $18 > 13$). Wynika to z "sztywnych" ograniczeń w poszczególnych komórkach tabeli (np. w p1 na Zmianie 2 musimy dać aż 7 aut, co podbija sumę).
\end{itemize}

% --- KONIEC ZADANIA 5 ---

% --- POCZĄTEK ZADANIA 6 ---
\section{Zadanie 6: Rozmieszczenie kamer (Set Cover)}

\subsection{Opis Modelu}
Problem optymalnego rozmieszczenia kamer monitoringu jest wariantem problemu pokrycia zbioru (Set Cover Problem). Celem jest minimalizacja liczby kamer potrzebnych do monitorowania wszystkich kontenerów.

\begin{itemize}
    \item $R = \{1, \dots, m\}$ - zbiór wierszy
    \item $C = \{1, \dots, n\}$ - zbiór kolumn
    \item $Q \subset R \times C$ - zbiór współrzędnych, w których znajdują się kontenery
    \item $k$ - zasięg widzenia kamery
\end{itemize}

\noindent \textbf{Zmienne decyzyjne:}
$$ x_{ij} \in \{0, 1\} \quad \forall i \in R, j \in C $$
gdzie $x_{ij} = 1$, jeśli w kwadracie $(i, j)$ umieszczono kamerę.

\noindent \textbf{Cel - minimalizacja liczby kamer:}
$$ \min \sum_{i=1}^{m} \sum_{j=1}^{n} x_{ij} $$

\noindent \textbf{Ograniczenia:}
\begin{enumerate}
    \item \textbf{Zakaz umieszczania na kontenerach:} Kamera nie może stać w miejscu zajętym przez kontener.
    $$ x_{ij} = 0 \quad \forall (i, j) \in Q $$
    
    \item \textbf{Pokrycie monitoringu:} Każdy kontener musi być obserwowany przez co najmniej jedną kamerę.
    $$ \sum_{(i, j) \in V_{(r,c)}} x_{ij} \ge 1 \quad \forall (r, c) \in Q $$
\end{enumerate}

\noindent \textbf{Definicja zbioru widoczności $V_{(r,c)}$:}
Zdefiniowano dwa warianty kształtu pola widzenia kamery:

\textbf{A. Kształt "Plus" (+):} Kamera widzi tylko wzdłuż osi (zgodnie z podstawową treścią zadania).
$$ V^{Plus}_{(r,c)} = \{ (i,j) : (i=r \land |j-c| \le k) \lor (j=c \land |i-r| \le k) \} $$

\textbf{B. Kształt "Kwadrat" ($\blacksquare$):} Kamera widzi cały obszar $(2k+1) \times (2k+1)$ (rozszerzenie eksperymentalne).
$$ V^{Kwadrat}_{(r,c)} = \{ (i,j) : |i-r| \le k \land |j-c| \le k \} $$

\subsection{Wyniki eksperymentu}
Wygenerowano instancję problemu o wymiarach $6 \times 6$ z 6 losowo rozmieszczonymi kontenerami. Przetestowano cztery scenariusze (kombinacje kształtu i zasięgu). Poniżej przedstawiono bezpośredni wydruk z programu:

\bigskip
% Używamy minipage, żeby całość wydruku była na jednej stronie (nierozrywalna)
\noindent
\begin{minipage}{\linewidth}
\begin{verbatim}
DANE: Teren 6x6, Liczba kontenerów: 6

==================================================
   ROZWIĄZYWANIE: KSZTAŁT = PLUS, ZASIĘG k = 1
==================================================
Minimalna liczba kamer: 6
Mapa (PLUS, k=1):
   1 2 3 4 5 6
1: C . . . . . 
2: K . . K C .
3: . K C . . .
4: . . . . K C
5: K C . . . .
6: . . K C . .
Lista kamer: (2, 1), (2, 4), (3, 2), (4, 5), (5, 1), (6, 3)

==================================================
   ROZWIĄZYWANIE: KSZTAŁT = PLUS, ZASIĘG k = 2
==================================================
Minimalna liczba kamer: 3
Mapa (PLUS, k=2):
   1 2 3 4 5 6
1: C . . . . .
2: . . . . C .
3: K . C . . .
4: . . . . K C
5: . C . . . .
6: . K . C . .
Lista kamer: (3, 1), (4, 5), (6, 2)

==================================================
   ROZWIĄZYWANIE: KSZTAŁT = KWADRAT, ZASIĘG k = 1
==================================================
Minimalna liczba kamer: 3
Mapa (KWADRAT, k=1):
   1 2 3 4 5 6
1: C . . . . .
2: . K . . C .
3: . . C . K .
4: . . . . . C
5: . C K . . .
6: . . . C . .
Lista kamer: (2, 2), (3, 5), (5, 3)

==================================================
   ROZWIĄZYWANIE: KSZTAŁT = KWADRAT, ZASIĘG k = 2
==================================================
Minimalna liczba kamer: 2
Mapa (KWADRAT, k=2):
   1 2 3 4 5 6
1: C . . . . .
2: K . . . C .
3: . . C . . .
4: . . . K . C
5: . C . . . .
6: . . . C . .
Lista kamer: (2, 1), (4, 4)
\end{verbatim}
\end{minipage}

\subsection{Analiza wyników}
\begin{itemize}
    \item \textbf{Wpływ kształtu:} Zmiana pola widzenia z "Plusa" na "Kwadrat" drastycznie zwiększa efektywność monitoringu. Dla zasięgu $k=1$, "Kwadrat" pozwolił zredukować liczbę kamer o połowę (z 6 do 3). Wynika to z faktu, że w metryce kwadratowej kamera pokrywa $(2k+1)^2$ pól, podczas gdy w "Plusie" tylko $4k+1$ pól.
    
    \item \textbf{Wpływ zasięgu ($k$):} Zwiększenie zasięgu z 1 do 2 również pozwala na znaczną redukcję floty kamer (w przypadku Plusa z 6 do 3).
    
    \item \textbf{Najlepsze rozwiązanie:} Najbardziej efektywny okazał się wariant \textbf{Kwadrat z $k=2$}, gdzie zaledwie **2 kamery** wystarczyły do pokrycia całego zestawu 6 kontenerów rozproszonych na planszy.
\end{itemize}

% --- KONIEC ZADANIA 6 ---

\end{document}

%%%%%%%%%%%%%%%%%%%%%%%%%%%%%%%%%%%%%%%%%%%%%%%%%%%%%%%%%%%%%%%%%%%%%%%%%%%%%%%
%
%   KONIEC PLIKU
%
%%%%%%%%%%%%%%%%%%%%%%%%%%%%%%%%%%%%%%%%%%%%%%%%%%%%%%%%%%%%%%%%%%%%%%%%%%%%%%%